\documentclass[10pt,a4paper]{amsart}
\usepackage{amsmath,amssymb,amsthm,graphicx}
% Retained minimal preamble for the Stacks Project chapter flat.tex (More on Flatness).
% Verbatim excerpt of the chapter preamble: Xy-pic setup (native preamble lines 14-18) and the native
% theorem declarations with the mathematical macros (native preamble lines 154-217). Copied mechanically;
% nothing added, removed or redefined. The wrapper class and the amsmath/amssymb/amsthm/graphicx base
% packages are supplied by the assembler, so the class/externaldocument part of the native preamble is omitted.
% For commutative diagrams we use Xy-pic
\usepackage[all]{xy}

% We use 2cell for 2-commutative diagrams.
\xyoption{2cell}
% Theorem environments.
%
\theoremstyle{plain}
\newtheorem{theorem}[subsection]{Theorem}
\newtheorem{proposition}[subsection]{Proposition}
\newtheorem{lemma}[subsection]{Lemma}

\theoremstyle{definition}
\newtheorem{definition}[subsection]{Definition}
\newtheorem{example}[subsection]{Example}
\newtheorem{exercise}[subsection]{Exercise}
\newtheorem{situation}[subsection]{Situation}

\theoremstyle{remark}
\newtheorem{remark}[subsection]{Remark}
\newtheorem{remarks}[subsection]{Remarks}

\numberwithin{equation}{subsection}

% Macros
%
\def\lim{\mathop{\mathrm{lim}}\nolimits}
\def\colim{\mathop{\mathrm{colim}}\nolimits}
\def\Spec{\mathop{\mathrm{Spec}}}
\def\Hom{\mathop{\mathrm{Hom}}\nolimits}
\def\Ext{\mathop{\mathrm{Ext}}\nolimits}
\def\SheafHom{\mathop{\mathcal{H}\!\mathit{om}}\nolimits}
\def\SheafExt{\mathop{\mathcal{E}\!\mathit{xt}}\nolimits}
\def\Sch{\mathit{Sch}}
\def\Mor{\mathop{\mathrm{Mor}}\nolimits}
\def\Ob{\mathop{\mathrm{Ob}}\nolimits}
\def\Sh{\mathop{\mathit{Sh}}\nolimits}
\def\NL{\mathop{N\!L}\nolimits}
\def\CH{\mathop{\mathrm{CH}}\nolimits}
\def\proetale{{pro\text{-}\acute{e}tale}}
\def\etale{{\acute{e}tale}}
\def\QCoh{\mathit{QCoh}}
\def\Ker{\mathop{\mathrm{Ker}}}
\def\Im{\mathop{\mathrm{Im}}}
\def\Coker{\mathop{\mathrm{Coker}}}
\def\Coim{\mathop{\mathrm{Coim}}}

% Boxtimes
%
\DeclareMathSymbol{\boxtimes}{\mathbin}{AMSa}{"02}

%
% Macros for moduli stacks/spaces
%
\def\QCohstack{\mathcal{QC}\!\mathit{oh}}
\def\Cohstack{\mathcal{C}\!\mathit{oh}}
\def\Spacesstack{\mathcal{S}\!\mathit{paces}}
\def\Quotfunctor{\mathrm{Quot}}
\def\Hilbfunctor{\mathrm{Hilb}}
\def\Curvesstack{\mathcal{C}\!\mathit{urves}}
\def\Polarizedstack{\mathcal{P}\!\mathit{olarized}}
\def\Complexesstack{\mathcal{C}\!\mathit{omplexes}}
% \Pic is the operator that assigns to X its picard group, usage \Pic(X)
% \Picardstack_{X/B} denotes the Picard stack of X over B
% \Picardfunctor_{X/B} denotes the Picard functor of X over B
\def\Pic{\mathop{\mathrm{Pic}}\nolimits}
\def\Picardstack{\mathcal{P}\!\mathit{ic}}
\def\Picardfunctor{\mathrm{Pic}}
\def\Deformationcategory{\mathcal{D}\!\mathit{ef}}

% --- appended 2026-10-12 (Chief): the project own definitions of the reference/slogan/history
% display environments, copied mechanically from the shared Stacks preamble lines 8-12
% (\usepackage{verbatim} plus the three \newenvironment lines, which the chapter preamble keeps
% verbatim from the project preamble). Nothing added, removed or redefined: these environments
% carry no mathematical content and the project itself typesets them as comments. Without these
% three lines a counted statement whose own text uses them fails with Environment slogan undefined.
\usepackage{verbatim}
\newenvironment{reference}{\comment}{\endcomment}
\newenvironment{slogan}{\comment}{\endcomment}
\newenvironment{history}{\comment}{\endcomment}
\usepackage{cleveref}
\crefname{originallemma}{original lemma}{original lemmas}
\Crefname{originallemma}{Original lemma}{Original lemmas}
\makeatletter
\newlabel{algebra-lemma-ascend-properties-modules}{{05A3}{source}{}{sourcepointer-1}{}}
\newlabel{algebra-lemma-ascend-properties-modules@cref}{{[originallemma][0][]05A3}{[1][0][]source}{}{}{}}
\makeatother
\crefname{originallemma}{original lemma}{original lemmas}
\Crefname{originallemma}{Original lemma}{Original lemmas}
\makeatletter
\newlabel{lemma-fibres-irreducible-flat-projective}{{05FS}{source}{}{sourcepointer-2}{}}
\newlabel{lemma-fibres-irreducible-flat-projective@cref}{{[originallemma][0][]05FS}{[1][0][]source}{}{}{}}
\makeatother
\crefname{originallemma}{original lemma}{original lemmas}
\Crefname{originallemma}{Original lemma}{Original lemmas}
\makeatletter
\newlabel{more-morphisms-lemma-Noetherian-approximation-combine}{{05FL}{source}{}{sourcepointer-3}{}}
\newlabel{more-morphisms-lemma-Noetherian-approximation-combine@cref}{{[originallemma][0][]05FL}{[1][0][]source}{}{}{}}
\makeatother
\crefname{originallemma}{original lemma}{original lemmas}
\Crefname{originallemma}{Original lemma}{Original lemmas}
\makeatletter
\newlabel{more-morphisms-lemma-Noetherian-approximation-flat}{{05FE}{source}{}{sourcepointer-4}{}}
\newlabel{more-morphisms-lemma-Noetherian-approximation-flat@cref}{{[originallemma][0][]05FE}{[1][0][]source}{}{}{}}
\makeatother
\crefname{originallemma}{original lemma}{original lemmas}
\Crefname{originallemma}{Original lemma}{Original lemmas}
\makeatletter
\newlabel{more-morphisms-lemma-Noetherian-approximation-geometrically-irreducible}{{05FH}{source}{}{sourcepointer-5}{}}
\newlabel{more-morphisms-lemma-Noetherian-approximation-geometrically-irreducible@cref}{{[originallemma][0][]05FH}{[1][0][]source}{}{}{}}
\makeatother
\crefname{originallemma}{original lemma}{original lemmas}
\Crefname{originallemma}{Original lemma}{Original lemmas}
\makeatletter
\newlabel{more-morphisms-lemma-Noetherian-approximation-geometrically-reduced}{{05FG}{source}{}{sourcepointer-6}{}}
\newlabel{more-morphisms-lemma-Noetherian-approximation-geometrically-reduced@cref}{{[originallemma][0][]05FG}{[1][0][]source}{}{}{}}
\makeatother
\title{A flat finitely presented algebra with geometrically integral fibres is projective over the base, without Noetherian hypotheses}
\author{The Stacks Project Authors}
\date{stacks-project:flat.tex:local-snapshot-20261011; GNU FDL 1.2 or later}
\begin{document}
\maketitle
\noindent\textit{Editorial note.} All delivered mathematical text is the author's own, copied from the pinned Stacks Project chapter file flat.tex; no statement or proof step was rewritten, reordered or removed.\par\medskip
\noindent\textit{Editorial note.} Preamble provenance: the chapter's own retained minimal preamble (sources/stacks-flat/retained-preamble.tex, lines 1-85) is a verbatim excerpt of the native Xy-pic setup and of the native theorem declarations and mathematical macros, followed by the project's own reference/slogan/history display-environment definitions copied mechanically from the shared preamble; the wrapper class and the amsmath/amssymb/amsthm/graphicx base packages are supplied by the assembler, and extra\_packages is left empty because everything needed is already in the retained preamble. No macro was invented or redefined.\par\medskip
\noindent\textit{Editorial note.} Same-chapter pointers are located at their \textbackslash{}label lines inside the pinned chapter file; cross-chapter pointers are located in the project's own tag list staged in the snapshot (sources/stacks-chow/tags.txt), at the verbatim lines carrying their canonical tags.\par\medskip
\noindent\textit{Allowed standard prerequisites (editorial).}\begin{itemize}
\item Approximation of the situation by one of finite type over the integers. More on Morphisms, labels more-morphisms-lemma-Noetherian-approximation-flat, more-morphisms-lemma-Noetherian-approximation-geometrically-reduced, more-morphisms-lemma-Noetherian-approximation-geometrically-irreducible and more-morphisms-lemma-Noetherian-approximation-combine (proof lines 2270-2273). Produces a cocartesian diagram with a base of finite type over the integers over which the same hypotheses hold, so that the Noetherian statement applies.
\item The Noetherian version of the projectivity statement, and ascent of projectivity along a base change. Same chapter flat.tex, label lemma-fibres-irreducible-flat-projective (proof line 2275, claimed as S14-algebra-fibres-irreducible-projective); Algebra, label algebra-lemma-ascend-properties-modules (proof line 2278). The first gives projectivity over the approximating base and the second ascends projectivity to the given base.
\end{itemize}\medskip
\ifdefined\hypertarget\hypertarget{sourcepointer-1}{}\fi
\noindent\textit{Source reference.} Original lemma 05A3: sources/stacks-chow/tags.txt, source lines 6836--6836. Navigation pointer for \textbackslash{}ref\{algebra-lemma-ascend-properties-modules\} used inside the counted span (line 2278); the argument that consumes it is retained in the span.\par\medskip
\ifdefined\hypertarget\hypertarget{sourcepointer-2}{}\fi
\noindent\textit{Source reference.} Original lemma 05FS: sources/stacks-flat/pinned-original-source.tex, source lines 2179--2179. Navigation pointer for \textbackslash{}ref\{lemma-fibres-irreducible-flat-projective\} used inside the counted span (line 2275); the argument that consumes it is retained in the span.\par\medskip
\ifdefined\hypertarget\hypertarget{sourcepointer-3}{}\fi
\noindent\textit{Source reference.} Original lemma 05FL: sources/stacks-chow/tags.txt, source lines 7030--7030. Navigation pointer for \textbackslash{}ref\{more-morphisms-lemma-Noetherian-approximation-combine\} used inside the counted span (line 2273); the argument that consumes it is retained in the span.\par\medskip
\ifdefined\hypertarget\hypertarget{sourcepointer-4}{}\fi
\noindent\textit{Source reference.} Original lemma 05FE: sources/stacks-chow/tags.txt, source lines 7023--7023. Navigation pointer for \textbackslash{}ref\{more-morphisms-lemma-Noetherian-approximation-flat\} used inside the counted span (line 2270); the argument that consumes it is retained in the span.\par\medskip
\ifdefined\hypertarget\hypertarget{sourcepointer-5}{}\fi
\noindent\textit{Source reference.} Original lemma 05FH: sources/stacks-chow/tags.txt, source lines 7026--7026. Navigation pointer for \textbackslash{}ref\{more-morphisms-lemma-Noetherian-approximation-geometrically-irreducible\} used inside the counted span (line 2272); the argument that consumes it is retained in the span.\par\medskip
\ifdefined\hypertarget\hypertarget{sourcepointer-6}{}\fi
\noindent\textit{Source reference.} Original lemma 05FG: sources/stacks-chow/tags.txt, source lines 7025--7025. Navigation pointer for \textbackslash{}ref\{more-morphisms-lemma-Noetherian-approximation-geometrically-reduced\} used inside the counted span (line 2271); the argument that consumes it is retained in the span.\par\medskip
\begin{lemma}
\label{lemma-fibres-irreducible-flat-projective-nonnoetherian}
Let $R$ be a ring. Let $R \to S$ be a ring map.
Assume
\begin{enumerate}
\item $R \to S$ is of finite presentation and flat, and
\item every fibre ring $S \otimes_R \kappa(\mathfrak p)$ is
geometrically integral over $\kappa(\mathfrak p)$.
\end{enumerate}
Then $S$ is projective as an $R$-module.
\end{lemma}

\begin{proof}
We can find a cocartesian diagram of rings
$$
\xymatrix{
S_0 \ar[r] & S \\
R_0 \ar[u] \ar[r] & R \ar[u]
}
$$
such that $R_0$ is of finite type over $\mathbf{Z}$, the map
$R_0 \to S_0$ is of finite type and flat with geometrically integral
fibres, see
More on Morphisms,
Lemmas \ref{more-morphisms-lemma-Noetherian-approximation-flat},
\ref{more-morphisms-lemma-Noetherian-approximation-geometrically-reduced},
\ref{more-morphisms-lemma-Noetherian-approximation-geometrically-irreducible},
and \ref{more-morphisms-lemma-Noetherian-approximation-combine}.
By
Lemma \ref{lemma-fibres-irreducible-flat-projective}
we see that $S_0$ is a projective $R_0$-module. Hence $S = S_0 \otimes_{R_0} R$
is a projective $R$-module, see
Algebra, Lemma \ref{algebra-lemma-ascend-properties-modules}.
\end{proof}

\end{document}
